%% §5 — the facilitation theorems. Drafted in Phase 1, in lockstep
%% with mechanization/Calculus/Frontier.lean; Lean names are cited
%% inline, and a statement without a Lean name says where its content
%% lives instead. Sufficiency theorems first, necessity ablations
%% after, assumptions typeset as the theorems' TCB.

\section{Why this explores the frontier}
\label{sec:theorems}

An explorer of a frontier must satisfy six properties or it is not an
explorer: it must draw a map its untrusted explorer cannot falsify
(F1), never lose ground (F2), know at every open point what would
extend the map there (F3), eventually reach every point inside the
true frontier (F4), know when a survey is finished (F5), and work in
any territory (F6). This section states each as a theorem over the
model of \Cref{sec:problem}, with its proof burden placed honestly:
mechanized, reduced to a named assumption, or empirical. The model is
small on purpose --- a registry is a set of items, a candidate is a
list of items, and answerability is the filtration of
\Cref{def:answerable}; its two monotonicities (stages shrink along
conditions, grow along registries) carry F2--F5.

\begin{theorem}[F1, unfalsifiable map]
\label{thm:sound}
Whatever the generator of pairs does --- blind or adversarial ---
the map's labels are truthful: (i) every existential answer is true
at the source, assuming only adequacy of the source interpreter; (ii)
every universal answer carries exactly the assurance class,
direction, and trusted base computed by the componentwise meet of its
route's contracts --- labels never overstate, because meets never
exceed their arguments.
\end{theorem}

\noindent
This is the instrument's own theory, consumed as an interface: (i) is
witness replay (\textsf{existential\_self\_certifying}, axiom-free),
(ii) is the contract algebra
(\textsf{weakest\_link\_universal}, \textsf{Contract.comp\_glb},
\textsf{lax\_universal\_transfer})~\citep{hurdygurdy2026calculus}.
The residue is stated, not hidden: a false universal in the corner
where no branch corroborates and no anchor covers is not excluded by
theorem --- it is bounded empirically (measured escape rates,
negative controls) and shrunk by the trust loop, and the map records
per answer which corner it sits in.

\begin{theorem}[F2, monotone exploration]
\label{thm:monotone}
If $G \subseteq G'$ then every question answerable at $G$ is
answerable at $G'$
{\normalfont(\textsf{Frontier.answerable\_mono}, axiom-free)}; and,
by the pair-level ratchet
{\normalfont(\textsf{ratchet\_preserves\_faithful})}, every standing
verdict stands.
\end{theorem}

\begin{theorem}[F3, complete local gradient]
\label{thm:gradient}
For every question unanswerable at $G$: a first failing condition
exists {\normalfont(\textsf{diagnosis\_total})}, is unique
{\normalfont(\textsf{diagnosis\_unique})}, never regresses under
growth {\normalfont(\textsf{diagnosis\_progress})}, and strictly
advances under any \emph{adequate} extension --- one after which some
admitted candidate survives one condition further
{\normalfont(\textsf{diagnosis\_strict\_progress})}.
\end{theorem}

\noindent
Totality is where the model earns the word \emph{complete}: the five
conditions are exhaustive \emph{by construction}, because
answerability is defined as their conjunction
(\Cref{def:answerable}), so the demand taxonomy is a theorem input,
not a design choice. What the platform's diagnosis \emph{names} ---
that the returned generation target, if admitted, discharges the
failing condition --- is the specification of \texttt{why\_not}, met
by construction in the implementation; it enters the theorems only as
the adequacy hypothesis. The one classical step
(\textsf{diagnosis\_total}) is itself informative: an unanswerable
question does not name its obstacle constructively --- the platform
computes it because its five conditions are decidable.

\begin{theorem}[F4, the chain lemma]
\label{thm:chain}
Along any growing chain of registries in which every diagnosis is met
by an adequate extension, $N$ steps answer the question
{\normalfont(\textsf{adequate\_chain\_answerable})}: the diagnosis
index is bounded by $N$ and strictly increases, so it runs out of
room.
\end{theorem}

\begin{assumption}[Fair enumeration]
\label{ass:fair}
Every demanded target is eventually attempted: the books aggregate
demand, and a human --- or a scoped mandate --- eventually acts on
standing demand.
\end{assumption}

\begin{assumption}[Gate liveness]
\label{ass:gate}
A sound design with adequate evidence is eventually admitted: builders
eventually produce, and the gate does not reject forever, a pair
discharging a named requirement that some sound pair can discharge.
\end{assumption}

\begin{theorem}[F4, relative completeness]
\label{thm:complete}
Under \Cref{ass:fair,ass:gate}, every question answerable at some
finite extension of $G$ over the candidate universe is eventually
answered by the loop.
\end{theorem}

\noindent
The mathematical content is \Cref{thm:chain}; the assumptions supply
its chain, and they are the theorem's honest trusted base --- both
are what the automation pipeline exists to provide, and both are
measured (queue throughput, gate escape rates) rather than provable.
The framing is deliberately that of relative
completeness~\citep{cook1978soundness}: complete relative to the
candidate universe and the domain's supply, not an oracle claim ---
undecidability, the adequacy floor of the source semantics, and the
finite anchor supply stand at every registry size.

\begin{theorem}[F5, saturation is a decidable, reachable fixpoint]
\label{thm:saturation}
For a pinned finite benchmark and a finite known set, the saturation
predicate of \Cref{def:saturation} is computable from the books ---
an emptiness check on the derived tier-2 set --- and the loop
restricted to in-set targets reaches it in finitely many iterations:
target signatures are finitely many, an admitted signature cannot
recur (dedup plus the ratchet), and each iteration retires at least
one.
\end{theorem}

\noindent
The finite-pool argument --- admitted signatures never recur, and
form a duplicate-free subset of the pool, so admissions run out of
room --- is mechanized as
{\normalfont\textsf{saturation\_terminates}}. The decidability half
is the implementation's contract: saturation is an exit code, not a
judgment.

\begin{theorem}[Currency: the status ratchet; conditional plans]
\label{thm:currency}
(i) Along any lifecycle a pair's tier only advances and its evidence
payload only accumulates
{\normalfont(\textsf{lifecycle\_ratchet}, axiom-free)}: a built pair
still carries the demand that created it. (ii) A mixed route's
conditional contract (\Cref{def:conditional}) is a lower bound on its
realized contract once every frontier hop is discharged by a pair
whose achieved contract dominates its requirement
{\normalfont(\textsf{conditional\_plan\_sound}, via monotonicity of
the meet)}: conditional plans never overpromise.
\end{theorem}

\begin{theorem}[F6, domain-genericity]
\label{thm:generic}
F1--F5 are stated and proved over the domain signature alone ---
abstract languages, items, conditions, contracts --- so they transfer
to any domain by instantiation. What does not transfer for free is
the supply side, and \Cref{sec:kit} states it as the explicit
obligation: the theorems are conditional on the kit, and the books
measure the condition.
\end{theorem}

\paragraph{Necessity: the ablations}
Each load-bearing feature of the instrument is needed for a specific
theorem, and dropping it breaks that theorem --- which is the precise
sense in which the current platform is the means to this paper's end.
Without \emph{determinism}, a square check cannot distinguish defect
from coin flip and F1(i)'s replay proves nothing. Without the
\emph{ratchet}, F2 fails and the curve can regress; a map that can
silently shrink is not a map. Without \emph{typed aborts} and the
fixed diagnosis order, the filtration of \Cref{def:answerable} has no
well-defined first failure and F3 degenerates to blind enumeration.
Without the \emph{existential/universal asymmetry} --- replay for
witnesses, graded meets for universals --- an untrusted generator
forges map entries and F1 falls adversarially. Without \emph{pinned
ingestion}, ``the benchmark'' is not a fixed set and F5's fixpoint is
not defined. And without the \emph{human valve} (or its mandate
generalization, which preserves it), \Cref{ass:fair}'s enumeration
has no scope and the books no meaning: demand-driven growth
presupposes someone accountable for what counts as demand.

\paragraph{The trusted base, collected}
F1(i) assumes source-interpreter adequacy; F1(ii)'s corner is
empirical; F4 assumes \Cref{ass:fair,ass:gate}; F5 assumes the pin
and the finiteness of the known set; F6 defers the supply side to the
kit. Nothing else is assumed, and each assumption is either measured
by the platform's own books or discharged by construction in the
implementation --- the reader can locate every claim's burden, which
is the property the whole architecture is built to have.
